% Options for packages loaded elsewhere
% Options for packages loaded elsewhere
\PassOptionsToPackage{unicode}{hyperref}
\PassOptionsToPackage{hyphens}{url}
\PassOptionsToPackage{dvipsnames,svgnames,x11names}{xcolor}
%
\documentclass[
  chinese,
  letterpaper,
  DIV=11,
  numbers=noendperiod]{scrartcl}
\usepackage{xcolor}
\usepackage[margin=2.5cm]{geometry}
\usepackage{amsmath,amssymb}
\setcounter{secnumdepth}{-\maxdimen} % remove section numbering
\usepackage{iftex}
\ifPDFTeX
  \usepackage[T1]{fontenc}
  \usepackage[utf8]{inputenc}
  \usepackage{textcomp} % provide euro and other symbols
\else % if luatex or xetex
  \usepackage{unicode-math} % this also loads fontspec
  \defaultfontfeatures{Scale=MatchLowercase}
  \defaultfontfeatures[\rmfamily]{Ligatures=TeX,Scale=1}
\fi
\usepackage{lmodern}
\ifPDFTeX\else
  % xetex/luatex font selection
\fi
% Use upquote if available, for straight quotes in verbatim environments
\IfFileExists{upquote.sty}{\usepackage{upquote}}{}
\IfFileExists{microtype.sty}{% use microtype if available
  \usepackage[]{microtype}
  \UseMicrotypeSet[protrusion]{basicmath} % disable protrusion for tt fonts
}{}
\makeatletter
\@ifundefined{KOMAClassName}{% if non-KOMA class
  \IfFileExists{parskip.sty}{%
    \usepackage{parskip}
  }{% else
    \setlength{\parindent}{0pt}
    \setlength{\parskip}{6pt plus 2pt minus 1pt}}
}{% if KOMA class
  \KOMAoptions{parskip=half}}
\makeatother
% Make \paragraph and \subparagraph free-standing
\makeatletter
\ifx\paragraph\undefined\else
  \let\oldparagraph\paragraph
  \renewcommand{\paragraph}{
    \@ifstar
      \xxxParagraphStar
      \xxxParagraphNoStar
  }
  \newcommand{\xxxParagraphStar}[1]{\oldparagraph*{#1}\mbox{}}
  \newcommand{\xxxParagraphNoStar}[1]{\oldparagraph{#1}\mbox{}}
\fi
\ifx\subparagraph\undefined\else
  \let\oldsubparagraph\subparagraph
  \renewcommand{\subparagraph}{
    \@ifstar
      \xxxSubParagraphStar
      \xxxSubParagraphNoStar
  }
  \newcommand{\xxxSubParagraphStar}[1]{\oldsubparagraph*{#1}\mbox{}}
  \newcommand{\xxxSubParagraphNoStar}[1]{\oldsubparagraph{#1}\mbox{}}
\fi
\makeatother


\usepackage{longtable,booktabs,array}
\usepackage{calc} % for calculating minipage widths
% Correct order of tables after \paragraph or \subparagraph
\usepackage{etoolbox}
\makeatletter
\patchcmd\longtable{\par}{\if@noskipsec\mbox{}\fi\par}{}{}
\makeatother
% Allow footnotes in longtable head/foot
\IfFileExists{footnotehyper.sty}{\usepackage{footnotehyper}}{\usepackage{footnote}}
\makesavenoteenv{longtable}
\usepackage{graphicx}
\makeatletter
\newsavebox\pandoc@box
\newcommand*\pandocbounded[1]{% scales image to fit in text height/width
  \sbox\pandoc@box{#1}%
  \Gscale@div\@tempa{\textheight}{\dimexpr\ht\pandoc@box+\dp\pandoc@box\relax}%
  \Gscale@div\@tempb{\linewidth}{\wd\pandoc@box}%
  \ifdim\@tempb\p@<\@tempa\p@\let\@tempa\@tempb\fi% select the smaller of both
  \ifdim\@tempa\p@<\p@\scalebox{\@tempa}{\usebox\pandoc@box}%
  \else\usebox{\pandoc@box}%
  \fi%
}
% Set default figure placement to htbp
\def\fps@figure{htbp}
\makeatother



\ifLuaTeX
\usepackage[bidi=basic,provide=*]{babel}
\else
\usepackage[bidi=default,provide=*]{babel}
\fi
% get rid of language-specific shorthands (see #6817):
\let\LanguageShortHands\languageshorthands
\def\languageshorthands#1{}


\setlength{\emergencystretch}{3em} % prevent overfull lines

\providecommand{\tightlist}{%
  \setlength{\itemsep}{0pt}\setlength{\parskip}{0pt}}



 


\usepackage{xeCJK}
% \setCJKmainfont{Source Han Serif SC} % 注释掉自定义字体，使用系统默认中文字体
\usepackage{tikz}
\usetikzlibrary{arrows.meta, calc, decorations.markings, positioning, shapes.geometric}
% Mermaid/PDF 图表支持
\usepackage{pgfplots}
\usepackage{graphicx}
% 改进的字体和排版支持
\usepackage{fontspec}
% PDF 超链接支持（若需要）
\usepackage{hyperref}
\KOMAoption{captions}{tableheading}
\makeatletter
\@ifpackageloaded{tcolorbox}{}{\usepackage[skins,breakable]{tcolorbox}}
\@ifpackageloaded{fontawesome5}{}{\usepackage{fontawesome5}}
\definecolor{quarto-callout-color}{HTML}{909090}
\definecolor{quarto-callout-note-color}{HTML}{0758E5}
\definecolor{quarto-callout-important-color}{HTML}{CC1914}
\definecolor{quarto-callout-warning-color}{HTML}{EB9113}
\definecolor{quarto-callout-tip-color}{HTML}{00A047}
\definecolor{quarto-callout-caution-color}{HTML}{FC5300}
\definecolor{quarto-callout-color-frame}{HTML}{acacac}
\definecolor{quarto-callout-note-color-frame}{HTML}{4582ec}
\definecolor{quarto-callout-important-color-frame}{HTML}{d9534f}
\definecolor{quarto-callout-warning-color-frame}{HTML}{f0ad4e}
\definecolor{quarto-callout-tip-color-frame}{HTML}{02b875}
\definecolor{quarto-callout-caution-color-frame}{HTML}{fd7e14}
\makeatother
\makeatletter
\@ifpackageloaded{caption}{}{\usepackage{caption}}
\AtBeginDocument{%
\ifdefined\contentsname
  \renewcommand*\contentsname{目录}
\else
  \newcommand\contentsname{目录}
\fi
\ifdefined\listfigurename
  \renewcommand*\listfigurename{图索引}
\else
  \newcommand\listfigurename{图索引}
\fi
\ifdefined\listtablename
  \renewcommand*\listtablename{表索引}
\else
  \newcommand\listtablename{表索引}
\fi
\ifdefined\figurename
  \renewcommand*\figurename{图}
\else
  \newcommand\figurename{图}
\fi
\ifdefined\tablename
  \renewcommand*\tablename{表}
\else
  \newcommand\tablename{表}
\fi
}
\@ifpackageloaded{float}{}{\usepackage{float}}
\floatstyle{ruled}
\@ifundefined{c@chapter}{\newfloat{codelisting}{h}{lop}}{\newfloat{codelisting}{h}{lop}[chapter]}
\floatname{codelisting}{列表}
\newcommand*\listoflistings{\listof{codelisting}{列表索引}}
\makeatother
\makeatletter
\makeatother
\makeatletter
\@ifpackageloaded{caption}{}{\usepackage{caption}}
\@ifpackageloaded{subcaption}{}{\usepackage{subcaption}}
\makeatother
\usepackage{bookmark}
\IfFileExists{xurl.sty}{\usepackage{xurl}}{} % add URL line breaks if available
\urlstyle{same}
\hypersetup{
  pdftitle={八字论命总纲（模板）},
  pdfauthor={你的名字},
  pdflang={zh-CN},
  colorlinks=true,
  linkcolor={blue},
  filecolor={Maroon},
  citecolor={Blue},
  urlcolor={Blue},
  pdfcreator={LaTeX via pandoc}}


\title{八字论命总纲（模板）}
\author{你的名字}
\date{}
\begin{document}
\maketitle


\section{最简七步（示例占位）}\label{ux6700ux7b80ux4e03ux6b65ux793aux4f8bux5360ux4f4d}

\begin{enumerate}
\def\labelenumi{\arabic{enumi}.}
\tightlist
\item
  立柱（定盘）\ldots\ldots{}
\item
  看日主旺衰（强/平/弱）\ldots\ldots{}
\item
  先调候（寒热燥湿）\ldots\ldots{}
\item
  看格局是否成而清\ldots\ldots{}
\item
  取用/定忌（主用/佐用/忌）\ldots\ldots{}
\item
  一验证（大运为底色，流年为按钮）\ldots\ldots{}
\item
  五行→行动（90天小清单）\ldots\ldots{}
\end{enumerate}

\begin{tcolorbox}[enhanced jigsaw, title=\textcolor{quarto-callout-note-color}{\faInfo}\hspace{0.5em}{注记}, leftrule=.75mm, toprule=.15mm, rightrule=.15mm, toptitle=1mm, colback=white, bottomtitle=1mm, breakable, colbacktitle=quarto-callout-note-color!10!white, coltitle=black, opacitybacktitle=0.6, titlerule=0mm, opacityback=0, arc=.35mm, left=2mm, bottomrule=.15mm, colframe=quarto-callout-note-color-frame]

\textbf{口诀}（可替换）：寒先温，热先凉，燥先润，湿先燥；强则泄制，弱则生扶；大运是底色，流年是按钮。

\end{tcolorbox}

\section{五行''生/克''关系图}\label{ux4e94ux884cux751fux514bux5173ux7cfbux56fe}

\begin{tikzpicture}[scale=1.2]
  % 颜色定义
  \definecolor{wuxingWood}{HTML}{4CAF50}
  \definecolor{wuxingFire}{HTML}{F44336}
  \definecolor{wuxingEarth}{HTML}{BCAAA4}
  \definecolor{wuxingMetal}{HTML}{9E9E9E}
  \definecolor{wuxingWater}{HTML}{2196F3}

  % 箭头样式
  \tikzset{
    >={Stealth[length=3mm,width=2mm]},
    sheng/.style={very thick, draw=black, ->},
    ke/.style={thick, dashed, draw=black, ->},
    elem/.style={circle, minimum size=12mm, inner sep=0pt, text=white, font=\bfseries}
  }

  % 五个顶点（顺时针：木、火、土、金、水）放在正五边形上
  \foreach \i/\name/\col in {
    90/木/wuxingWood,
    18/火/wuxingFire,
    -54/土/wuxingEarth,
    -126/金/wuxingMetal,
    162/水/wuxingWater
  }{
    \path ({cos(\i)},{sin(\i)}) node[elem, fill={\col}] (\name) {\name};
  }

  % 生（顺时针环）
  \draw[sheng] (木) edge[bend left=12] (火);
  \draw[sheng] (火) edge[bend left=12] (土);
  \draw[sheng] (土) edge[bend left=12] (金);
  \draw[sheng] (金) edge[bend left=12] (水);
  \draw[sheng] (水) edge[bend left=12] (木);

  % 克（五角星）
  \draw[ke] (木) -- (土);
  \draw[ke] (土) -- (水);
  \draw[ke] (水) -- (火);
  \draw[ke] (火) -- (金);
  \draw[ke] (金) -- (木);
\end{tikzpicture}

\begin{tcolorbox}[enhanced jigsaw, title=\textcolor{quarto-callout-note-color}{\faInfo}\hspace{0.5em}{注记}, leftrule=.75mm, toprule=.15mm, rightrule=.15mm, toptitle=1mm, colback=white, bottomtitle=1mm, breakable, colbacktitle=quarto-callout-note-color!10!white, coltitle=black, opacitybacktitle=0.6, titlerule=0mm, opacityback=0, arc=.35mm, left=2mm, bottomrule=.15mm, colframe=quarto-callout-note-color-frame]

\textbf{说明}: 此图展示了五行相生和相克的关系，形成完美的圆形闭环布局：
- \textbf{相生循环}（实线箭头）：木生火，火生土，土生金，金生水，水生木
- \textbf{相克循环}（虚线箭头）：木克土，土克水，水克火，火克金，金克木

\textbf{备选方案}: 如果需要更高质量/可控配色的图表，可以： 1. 编译
\texttt{tikz/wuxing.tex} 生成PDF图表 2.
然后在此处插入：\texttt{!{[}五行生克图{]}(tikz/wuxing.pdf)\{width="\textbackslash{}\textbackslash{}textwidth"\}}

\end{tcolorbox}

\section{用神→行动映射（示例表）}\label{ux7528ux795eux884cux52a8ux6620ux5c04ux793aux4f8bux8868}

\begin{longtable}[]{@{}
  >{\raggedright\arraybackslash}p{(\linewidth - 4\tabcolsep) * \real{0.3333}}
  >{\raggedright\arraybackslash}p{(\linewidth - 4\tabcolsep) * \real{0.3333}}
  >{\raggedright\arraybackslash}p{(\linewidth - 4\tabcolsep) * \real{0.3333}}@{}}
\toprule\noalign{}
\begin{minipage}[b]{\linewidth}\raggedright
五行
\end{minipage} & \begin{minipage}[b]{\linewidth}\raggedright
功能倾向（可编辑）
\end{minipage} & \begin{minipage}[b]{\linewidth}\raggedright
行动例子（90天可执行）
\end{minipage} \\
\midrule\noalign{}
\endhead
\bottomrule\noalign{}
\endlastfoot
木 & 生长、学习、拓新 & 每周新增一个学习模块；日常拉伸；远景规划回顾 \\
火 & 表达、体温、见光、社交 &
每日晨光10--20min；每周一次分享/演讲；日记输出 \\
土 & 秩序、边界、收尾、稳定 &
看板收尾率≥80\%；拒绝越界任务；固定复盘时段 \\
金 & 规则、工具、裁剪、复盘 & 模板化常用流程；重构工具链；周复盘+删减 \\
水 & 睡眠、补水、专注、移动 & 7.5h 睡眠；2L 饮水；远距离步行/短途旅行 \\
\end{longtable}

\section{口袋口诀（示例占位）}\label{ux53e3ux888bux53e3ux8bc0ux793aux4f8bux5360ux4f4d}

\begin{itemize}
\tightlist
\item
  年看背景，月看气，日看我，时看舞台。
\item
  成格从格优先，清纯胜过堆料。
\item
  神煞当形容词，不当判决书。
\item
  结论需可验证，并标注前提条件（时间协议/时差/节令）。
\end{itemize}




\end{document}
